\documentclass[11pt,a4paper]{article}
\usepackage[T1]{fontenc}
\usepackage[utf8]{inputenc}
\usepackage{lmodern,amsmath,amssymb,booktabs,listings,hyperref,microtype}
\usepackage[margin=26mm]{geometry}
\hypersetup{colorlinks=true,urlcolor=blue,linkcolor=blue,pdftitle={The Maximum Determinant of a 55 by 55 Binary Circulant Matrix}}
\lstset{basicstyle=\ttfamily\footnotesize,breaklines=true,columns=fullflexible,keepspaces=true}
\setlength{\parindent}{0pt}
\setlength{\parskip}{5pt}
\setlength{\emergencystretch}{2em}
\title{The Maximum Determinant of a\\55 $\times$ 55 Binary Circulant Matrix}
\author{Reproducible computational research note}
\date{9 September 2026}
\begin{document}
\maketitle
\section*{Abstract}

For a binary word a of length 55, let C(a) have entries a[(j\ensuremath{-}i) mod 55]. We determine the maximum of det C(a) over all \ensuremath{2^55} words:

\[
D_{01}(55)=134694094094758395331307111329132.
\]

All maximizing words have weight 28 and form exactly one orbit under the affine group of the cyclic group of order 55. The orbit has 2,200 elements and trivial stabilizer. The proof combines complement duality, integer periodic correlations, fourth-moment product bounds with a spectral cap, cyclotomic folds of orders 5 and 11, exact correlation-profile enumeration, and binary torus lifts. Only 6,845,230 ordered correlation profiles require explicit enumeration. They leave 2,316 unit-orbit representatives and 16,084 margin tasks. A meet-in-the-middle lift examines 15,487,882,832 joined words and returns two normalized words, belonging to the same affine class. A separate verifier recomputes the mathematical bounds and exact norms, regenerates the profile enumerations, and checks all lifts using a different column split. Every component is SHA-256-bound to the global certificate. Neither floating-point scores nor an uncompleted search are used as evidence of maximality.

\section*{1. Introduction and previous exact frontier}

Write D01(n) for the largest absolute determinant of an n \ensuremath{\times} n binary circulant matrix. Brent and Yedidia developed exact necklace enumeration and Fourier determinant algorithms; the updated version of their paper reports values through n=53 [1]. The OEIS A086432 b-file inspected on 2026-09-09 ends at n=53 [2]. ALETHEIA-54 supplies a computational proof at n=54 based on integer correlation profiles and folded binary lifts [3]. The present calculation addresses n=55, whose factorization 5\ensuremath{\cdot}11 replaces the parity fold available at even order.

\textbf{Theorem 1 (global maximum and complete classification).}

\[
\max_{a\in\{0,1\}^{55}}\det C(a)
=134694094094758395331307111329132=:M.
\]

A lexicographically minimal representative of the unique maximizing affine class is

\begin{lstlisting}
0000000100011111011011001101011011010100011110101000111
\end{lstlisting}

Its weight is 28, its affine stabilizer is trivial, and all 2,200 maximizing words are obtained by the transformations \ensuremath{a_i} \ensuremath{\mapsto} \ensuremath{a_{ui+t},} with u a unit modulo 55 and t any residue modulo 55.

The complete finite computations supporting the theorem are recorded in \path|certificates/order55_global/|. The scope is binary circulants, rather than all binary matrices of order 55 or unrestricted \ensuremath{\pm}1 matrices.

\section*{2. Odd-order Fourier formulation}

Let \ensuremath{f(x)=\sum _{t=0}^{54}a_t} \ensuremath{x^t,} \ensuremath{\zeta}=exp(2\ensuremath{\pi}i/55), and \ensuremath{k=\sum a_t.} The eigenvalues of C(a) are \ensuremath{f(\zeta ^j),} possibly in the reverse frequency order according to the eigenvector convention. For real a and odd order all nonzero frequencies pair with their conjugates. Put

\[
q_j=|f(\zeta^j)|^2,\quad 1\le j\le27.
\]

Then

\[
\det C(a)=k\prod_{j=1}^{27}q_j\ge0,\qquad
\sum_{j=1}^{27}q_j=\frac{55k-k^2}{2}.
\]

The nonnegativity makes absolute-value notation unnecessary for actual 55-bit words. Absolute products remain useful when screening formal correlation profiles which have not yet been proved positive semidefinite.

\section*{3. Complement duality and affine symmetry}

For 0<k<55, complement changes the zero-frequency eigenvalue from k to 55\ensuremath{-}k and negates all 54 other eigenvalues. Hence

\[
k\det C(1-a)=(55-k)\det C(a),
\qquad
D_{\rm circ}(55,55-k)=\frac{55-k}{k}D_{\rm circ}(55,k).
\]

Complement is a bijection between the two weight strata. Both constant words are singular. Since (55\ensuremath{-}k)/k>1 for k\ensuremath{\le}27,

\[
D_{01}(55)=\max_{1\le k\le27}\frac{55-k}{k}D_{\rm circ}(55,k).
\]

Translation multiplies Fourier eigenvalues by phases with product \ensuremath{\zeta ^{t\cdot 55\cdot 54/2}=1.} Multiplication of the indices by a unit permutes frequencies. Thus the affine group of order 55\ensuremath{\varphi}(55)=2,200 preserves the determinant. Reversal is its u=\ensuremath{-}1 element. We use actual stabilizers rather than assuming all orbits have size 2,200.

Correlations are invariant under translation and reversal identifies opposite shifts. The effective unit action on the 27 paired shifts has order 20. Its shift orbits have sizes 20, 5, and 2, according as the gcd with 55 is 1, 5, or 11.

\section*{4. Periodic correlations and weight reduction}

Define \ensuremath{c_s=\sum _t} \ensuremath{a_t} \ensuremath{a_{t+s}.} Then \ensuremath{c_0=k,} \ensuremath{c_s=c_{55- s},} and

\[
2\sum_{s=1}^{27}c_s=k(k-1),\qquad
k^4+2\sum_{j=1}^{27}q_j^2
=55\left(k^2+2\sum_{s=1}^{27}c_s^2\right).
\]

These follow by Fourier orthogonality, since the Fourier transform of the correlation sequence is \ensuremath{|f(\zeta ^j)|^2.} If T nonnegative integer units are allocated among h coordinates, write T=hb+r, 0\ensuremath{\le}r<h. The minimum square sum is (h\ensuremath{-}r)b\textsuperscript{2}+r(b+1)\textsuperscript{2}: transferring one unit between coordinates differing by at least two strictly decreases the square sum.

The first spectral moment gives the fixed-row-sum bound

\[
\det C(a)\le k\left(\frac{k(55-k)}{54}\right)^{27}.
\]

For the complementary high-weight word the leading factor is 55\ensuremath{-}k. Combining this bound with the integer fourth moment excludes k=1,...,24 at threshold M. The certificate lists each weight pair separately. Only k=25,26,27 remain, with the following necessary bounds:

\begin{center}
\footnotesize
\setlength{\tabcolsep}{4pt}
\begin{tabular}{lrr}
\toprule
Lower weight k & Minimum \ensuremath{\sum} \ensuremath{c_s^2} & Maximum \ensuremath{\sum} \ensuremath{c_s^2} to reach M \\
\midrule
25 & 3336 & 3338 \\
26 & 3913 & 3919 \\
27 & 4563 & 4570 \\
\bottomrule
\end{tabular}
\end{center}

All sums in this table use s=1,...,27. For k=27, \ensuremath{d_s=c_s- 13} has sum zero, so \ensuremath{E=\sum d_s^2} is even. Thus the last row permits only E\ensuremath{\le}6.

\section*{5. Exact fourth-moment product bounds}

Consider h nonnegative variables with sum S and square sum at least L. Contracting toward the common mean reduces the square sum and increases the product, by concavity of log and continuity at zero. It suffices to use square sum max(L,S\textsuperscript{2}/h). At a positive constrained stationary point,

\[
1/x=\alpha+2\beta x,
\]

so there are at most two distinct coordinates. If the smaller value occurs m times, let \ensuremath{\mu}=S/h and V=L\ensuremath{-}S\textsuperscript{2}/h. The two values are

\[
\mu-\sqrt{\frac{V(h-m)}{hm}},\qquad
\mu+\sqrt{\frac{Vm}{h(h-m)}}.
\]

Enumerate every m=1,...,h\ensuremath{-}1, retaining positive values. Boundary products with a zero coordinate vanish. Rational outward square-root intervals bound each stationary product from above. A zero-variance case is handled directly. The implementation uses integer square roots at a decimal scale of \ensuremath{10^50;} no floating-point comparison decides a pruning step.

\subsection*{5.1. A spectral cap from each correlation multiset}

The unbounded stationary problem is sometimes too permissive. Let b=floor(k(k\ensuremath{-}1)/54), and sort the 27 deviations \ensuremath{c_s- b.} For any nonzero frequency,

\[
q_j=k-b+\sum_{s=1}^{27}(c_s-b)\,2\cos(2\pi js/55).
\]

The rearrangement inequality bounds this expression by pairing the sorted deviations with sorted cosine coefficients. There are only three possible coefficient multisets, represented by j=1,5,11. The largest of their outward rational bounds is a uniform cap B for every \ensuremath{q_j.}

To maximize the product subject also to \ensuremath{x_i\le B,} enumerate r coordinates at the upper boundary B. The remaining h\ensuremath{-}r coordinates have sum S\ensuremath{-}rB and square sum L\ensuremath{-}rB\textsuperscript{2}; their interior stationary points again have at most two values. Discard infeasible sums, negative coordinates, and a larger root provably above B. Include all-zero-variance cases. Zero lower-boundary coordinates contribute product zero. This exhausts the compact feasible set and proves the capped bound used in the certificate.

Pi is enclosed using Machin's identity and rational alternating series for arctangent. Rational Taylor remainders then enclose the cosines. The native screen uses these intervals at dyadic scale \ensuremath{2^24.} Its multiply-and- round steps are upward, and the builder and verifier separately check integer inequalities that bound every unsigned-128-bit intermediate.

\section*{6. The 5 \ensuremath{\times} 11 cyclotomic decomposition}

The frequency sectors of root orders 1,5,11,55 have dimensions 1,4,10,40. Consequently

\[
\det C(a)=kN_5(a)N_{11}(a)N_{55}(a),\qquad
N_m(a)=\operatorname{Res}(f,\Phi_m).
\]

All three cyclotomic degrees are even, so exchanging the resultant operands has positive sign. Their factors pair conjugately and the norms are nonnegative. The paired sector dimensions are 2,5,20.

Let \ensuremath{r_i=\sum _{t\equiv i} (mod \ensuremath{5)}a_t} and \ensuremath{s_j=\sum _{t\equiv j} (mod \ensuremath{11)}a_t.} Then \ensuremath{0\le r_i\le 11} and \ensuremath{0\le s_j\le 5;} N5 depends only on r and N11 only on s. For either modulus m, write \ensuremath{h_i=\sum _{t\equiv i} (mod \ensuremath{m)}c_t.} This is the cyclic correlation of the folded margin vector. The sector moments are

\[
S_1^{(m)}=\frac{m\sum r_i^2-k^2}{2},\qquad
S_2^{(m)}=\frac{m\sum h_i^2-k^4}{2}.
\]

Subtract the two small sectors from the full moments for the primitive 55 sector. For fixed r,s its first moment is

\[
T=\frac{55k+k^2-5\sum r_i^2-11\sum s_j^2}{2},
\]

and its norm is at most \ensuremath{(T/20)^20.} The exploratory calculation also applied this combined bound and the separate sector fourth moments. The final proof uses an even simpler necessary bound for each individual fold: if its sector has h pairs and sum T, the remaining 27\ensuremath{-}h pairs have sum R=k(55\ensuremath{-}k)/2\ensuremath{-}T, and

\[
(55-k)N_m\left(\frac{R}{27-h}\right)^{27-h}\ge M
\]

is necessary. Both inequalities are rational. Every bounded integer fold passing the necessary condition is enumerated, up to its full affine symmetry, and its norm is calculated exactly in a prime field. A Parseval bound proves that each such small norm is below the prime, so there is no modular ambiguity. The independent verifier checks every retained norm again by an integer resultant.

\begin{center}
\footnotesize
\setlength{\tabcolsep}{4pt}
\begin{tabular}{lrr}
\toprule
k & Mod-5 affine representatives & Mod-11 affine representatives \\
\midrule
25 & 17 & 321 \\
26 & 27 & 560 \\
27 & 26 & 660 \\
\bottomrule
\end{tabular}
\end{center}

All unit orientations of their folded correlations are retained. Translation acts trivially on correlations. Independent translations of the two margin vectors lift to a translation of the original word by the Chinese remainder theorem. Therefore using the minimum rotation of each oriented margin loses no necessary binary lift.

\section*{7. Difference-set obstruction and winner structure}

The ideal weight-27 correlations would all equal 13, giving a cyclic (55,27,13) difference set. Here is a short algebraic nonexistence proof. Put \ensuremath{\alpha}=f(\ensuremath{\zeta}5). Ideal correlations imply \ensuremath{\alpha} conjugate(\ensuremath{\alpha})=14. Modulo 2, \ensuremath{\Phi}5 is irreducible since ord5(2)=4, so Z[\ensuremath{\zeta}5]/(2) is a field. The product is zero there; conjugation is an automorphism, so \ensuremath{\alpha} is divisible by 2. Writing \ensuremath{\alpha}=2\ensuremath{\beta} would make \ensuremath{\beta} conjugate(\ensuremath{\beta})=7/2, contrary to the fact that a rational algebraic integer is an integer. Published corroboration appears in Brai\'{c}, Theorem 4.4 and p.297, using p=2,w=5,j=2 and k\ensuremath{-}\ensuremath{\lambda}=14 [4].

The maximizing word is not a two-level almost difference set in the usual sense [5]. Its nonzero correlations have three values: 13 occurs four times, 14 occurs 46 times, and 15 occurs four times. In the displayed canonical representative, the independent exceptional shifts are 2,6 (value 15) and 7,15 (value 13). The complement has E=4 about correlation 13. Its associated \ensuremath{\pm}1 periodic autocorrelations are \ensuremath{-}5,\ensuremath{-}1,3, with total nonzero-shift squared energy 182. We do not claim a Bernasconi ground-state theorem or a new algebraic construction. The trivial affine stabilizer rules out a nontrivial multiplier symmetry even after translation.

The winner has

\begin{lstlisting}
mod-5 margin:  [5,4,8,6,5]
mod-11 margin: [3,1,4,2,3,1,1,3,3,3,4]
N5  = 131
N11 = 434039
N55 = 84603917386870862636041
\end{lstlisting}

Their product times 28 is exactly M.

\section*{8. Complete correlation-profile enumeration}

For each remaining k enumerate every sorted multiset of 27 integer correlations in [0,k] with sum k(k\ensuremath{-}1)/2 and permitted square sum. A balanced-square lower bound prunes only impossible recursion branches. The multinomial coefficient gives the exact number of ordered profiles. First apply the capped moment bound to entire multisets. For each surviving multiset enumerate every coordinate assignment.

A formal profile must have folded correlations matching retained margin vectors for both moduli. After this exact membership test, retain its lexicographically minimal representative under the effective 20-element unit group. Dyadic outward spectral products provide a safe first screen. Remaining products are evaluated exactly using two 61-bit Fourier fields and signed CRT. Even if a formal profile is not PSD, Cauchy--Schwarz and AM-GM give

\[
\left|(55-k)\prod q_j\right|
\le(55-k)\left(\frac{\sum q_j^2}{27}\right)^{27/2}.
\]

A squared integer comparison proves that signed CRT capacity exceeds this bound. Only profiles with exact absolute product at least M become lift targets. For any realizable profile this is its actual complementary determinant, so the screening condition covers equality as well as a possible improvement.

\begin{center}
\footnotesize
\setlength{\tabcolsep}{4pt}
\begin{tabular}{lrrrr}
\toprule
k & Allowed profiles & Enumerated & Canonical & Lift targets \\
\midrule
25 & 407,277 & 2,925 & 55 & 55 \\
26 & 32,178,654 & 816,102 & 22,375 & 281 \\
27 & 6,043,753 & 6,026,203 & 155,558 & 1,980 \\
Total & 38,629,684 & 6,845,230 & 177,988 & 2,316 \\
\bottomrule
\end{tabular}
\end{center}

The first column is an overapproximation of possible integer correlations, not a count of binary words. The explicitly enumerated columns are actual completed work, not projected throughput.

\section*{9. Exhaustive binary torus lifts}

Identify the torus coordinate (i,j)\ensuremath{\in}Z5\ensuremath{\times}Z11 with 11i+45j modulo 55. Each of its 11 columns is one of 32 five-bit masks. For a target correlation profile, enumerate every oriented pair of compatible margins and normalize independent row and column translations. This gives 16,084 tasks.

Split the columns into five and six. A half-assignment stores its binary word and a seven-coordinate signature: the five row counts and the two within-column correlation sums for shifts 11 and 22. The latter are additive across columns. The left half is sorted by its packed signature. For each right half, compute the exact complementary signature and join every matching left entry. No hash collision is used as equality: the packed integer signature is injective on the bounded coordinate ranges. Every joined word has the requested margins and these two correlations. Test every other independent shift by rotated 55-bit intersection and population count. A remaining match is a full binary lift.

This is a complete meet-in-the-middle enumeration of each margin fiber. Nonnegative partial row counts and within-column correlations exceeding their targets can safely be pruned. No floating-point operation is involved. The computation tests 15,487,882,832 joined words and finds exactly two normalized solutions. An independent audit repeats every task with a six-plus-five split and obtains identical joined-word counts and identical solution sets. A separate small fiber was compared with all 78,125 direct column assignments, using three different splits.

\section*{10. Global maximality and classification proof}

Any binary word with determinant at least M has a lower complementary weight k\ensuremath{\in}\{1,...,27\}. The exact weight screen excludes k\ensuremath{\le}24. For each remaining k, the moment equations place its correlation multiset in the complete enumerated list. Capped moment bounds exclude some multisets; the others are exhausted in their entirety. Its margins must pass both folded screens, and a unit transformation takes its paired correlations to one of the 2,316 retained representatives. Independent translations then take both oriented margins to recorded minimum rotations. The complete lift of that margin pair therefore contains a transformed copy of the word.

The only two lifted words have complementary determinant M. Their higher-weight complements belong to the same affine class. Exact matrix determinants, finite-field Fourier CRT, and cyclotomic resultant products all agree. Thus no word exceeds M, and every word attaining M is in the reported class. Enumerating its affine action gives 2,200 distinct words, so the stabilizer has order one. This proves Theorem 1.

\section*{11. Independent certification and reproducibility}

The global JSON file binds a hash manifest and all component files. \path|verify_order55_global.py| imports neither the search nor the production certificate builder. It independently recomputes the rational moment and cap formulas, verifies a separate Machin/Taylor enclosure, checks all retained norms by resultants, recompiles hash-identified C++ sources, and replays the folded and correlation enumerations. It derives every required margin task afresh and checks task identities, multiplicities, output words, exact determinants, and affine classification. Its optional full mode also replays every lift with a different split. Both audit modes and the complete full-mode replay passed in this run.

The finite-field primes are 2305843009213696591 and 2305843009213697141; each is prime and is 1 modulo 55. The winner residues are respectively 1234839479219971141 and 1202711640743500591. The builder and verifier check the applicable reconstruction and fixed-point intermediate bounds.

On the recorded Windows host, use:

\begin{lstlisting}
& ./.venv/Scripts/python.exe scripts/build_order55_global.py
& ./.venv/Scripts/python.exe scripts/verify_order55_global.py --full-lifts
& ./.venv/Scripts/python.exe -m pytest
& ./.venv/Scripts/python.exe -m pip check
\end{lstlisting}

The same Python entry points accept an ordinary Python interpreter with SymPy and the package installed. Native compilation needs Clang/GCC C++20, unsigned \path|__int128|, and C++20 population count. No upstream repository, stochastic search, GPU, network access, or proprietary solver is needed to reconstruct the final proof. The stored C++ build commands and source hashes identify the original computation. Paths and elapsed times may differ across hosts; the discrete coverage counts, words, profile products, and classification must agree.

\section*{12. Prior-work check}

On 2026-09-09 the check included the OEIS entry and b-file, Brent's maximal determinant pages, the latest Brent--Yedidia arXiv revision, ALETHEIA-54, and web queries targeting arXiv, GitHub, Zenodo, and the exact determinant integer. No prior exact order-55 determination was located. Search engines have incomplete coverage; this is not an absolute-priority claim.

We found no prior published exact determination in the sources checked
as of 2026-09-09.

The detailed retrieval log is in \path|references/order55/SOURCES.md|.

\section*{Appendix A. Previously established small-weight strata}

The frozen starting repository independently certified \ensuremath{D_{\rm circ}(55,k)} for k=1,...,7. Complement duality immediately supplies the paired high-weight strata and hence additional supporting bounds in the global problem. The present global screen does not need to rerun their large subset exhaustions. Their existing certificate and verifier are retained.

\begin{center}
\footnotesize
\setlength{\tabcolsep}{4pt}
\begin{tabular}{lr}
\toprule
k & \ensuremath{D_{\rm circ}(55,k)} \\
\midrule
1 & 1 \\
2 & 2048 \\
3 & 1564031349 \\
4 & 6424318816256 \\
5 & 12556202952818275 \\
6 & 4342461217049330166 \\
7 & 570999857161750276477 \\
\bottomrule
\end{tabular}
\end{center}

\section*{References}

\par\medskip\noindent 1. R. P. Brent and A. B. Yedidia, \emph{Computation of Maximal Determinants of Binary Circulant Matrices}, Journal of Integer Sequences 21 (2018), Article 18.5.6; updated arXiv:1801.00399v6, 20 February 2021. \url{https://arxiv.org/abs/1801.00399v6}
\par\medskip\noindent 2. OEIS Foundation, \emph{A086432: Maximum determinant of binary circulants}, entry and b-file inspected 9 September 2026. \url{https://oeis.org/A086432}, \url{https://oeis.org/A086432/b086432.txt}
\par\medskip\noindent 3. K. Terauchi, \emph{ALETHEIA-54}, upstream commit c1332afa5d406373ec47f63dc3f187aee4a6cecf, retrieved and checked 9 September 2026. Source MIT; documentation and proof data CC BY 4.0. \url{https://github.com/ruu413/aletheia-54}
\par\medskip\noindent 4. S. Brai\'{c}, \emph{Primitive symmetric designs with at most 255 points}, Glasnik Matematicki 45(65) (2010), 291--305, Theorem 4.4 and p.297. \url{https://web.math.pmf.unizg.hr/glasnik/45.2/45(2)-01.pdf} See also E. S. Lander, \emph{Symmetric Designs: An Algebraic Approach}, Cambridge University Press, 1983, pp.131,134, cited there.
\par\medskip\noindent 5. K. Nowak, \emph{A Survey on Almost Difference Sets}, arXiv:1409.0114v1, 30 August 2014. \url{https://arxiv.org/abs/1409.0114v1}
\end{document}
